\subsection{Consistencia en la clase débil}\label{sec:B5T05}
\begin{teorema}[Consistencia mesoscópica]\label{thm:meso}
Para $C,\eta,\gamma$ fijos existen $\varepsilon_0,k_0,C_*$ independientes
de $n/k$ con la siguiente propiedad. Si el padre satisface las condiciones de la Sección~\ref{sec:B5T01} con
$\varepsilon\le\varepsilon_0$, $k\ge k_0$, entonces todos sus hijos legales
tienen el punto estacionario pequeño construido en la Sección~\ref{sec:B5T03} y
\[
 \boxed{\left|\log\rho(B)\right|\le C_*/k^2,\qquad
 \rho(B)=\sum_{c\in r}\exp[L(B+c)-L(B)].}
\]
\end{teorema}
En el enunciado se fija una fila $r$ y se incluyen todos sus hijos; $L$ es exactamente
el predictor canónico: término de árbol, factor local, cofactor RC y covolumen.
No se identifica su marginal corregida con la marginal verdadera de las soluciones.
El predictor selecciona la rama pequeña: la hessiana cambiada de signo
es definida positiva en la región convexa $p<1/4$, con vacancias positivas.
Así, el punto de Newton es su único punto estacionario en esa rama, y no una selección nueva.
El desplazamiento total es $y=z+s+O_{\rm rel}(k^{-2})$, donde $s$ es el paso de Newton desde $\bar p$.
Se reserva margen en $\varepsilon_0,C,\eta,\gamma$ para todos los hijos.

\begin{proof}
\emph{Adaptación de la prueba macroscópica.}
La prueba sigue la de la Sección~\ref{sec:macro}; sus identidades algebraicas y su normalización se
conservan. Se sustituyen sus usos de $k\ge\delta_0n$ por la Sección~\ref{sec:B5T02} para la inversa de
líneas y la respuesta, la Sección~\ref{sec:B5T03} para el paso de Newton, la masa de la diferencia de
Schur y los restos cúbicos, y la Sección~\ref{sec:B5T04} para las trazas, sus cambios y la resolvente.
El reequilibrado utiliza sólo las $2k$ líneas RC.
En ningún punto se exige un número $O(k)$ de diagonales. La correspondencia es la siguiente.
\begin{center}
\renewcommand{\arraystretch}{1.25}
\begin{tabular}{@{}>{\raggedright\arraybackslash}p{0.52\textwidth}
                >{\raggedright\arraybackslash}p{0.43\textwidth}@{}}
\toprule
Paso de la prueba macroscópica & Lo sustituye\\
\midrule
$G=I+O_{\max}(k^{-1})$ y $\|G\|_\infty=O(1)$ (Sección~\ref{sec:B3T05}) &
 Sección~\ref{sec:B5T02}, complemento de Schur\\
$\|P^{-1}R\|_\infty=O(1)$ (Sección~\ref{sec:B3T06}) &
 Sección~\ref{sec:B5T02}, respuesta ponderada\\
Paso de Newton y resto cúbico de la acción (Sección~\ref{sec:B4T02}) &
 Sección~\ref{sec:B5T03}\\
$\|\mathcal E\|_F^2$, $\sum_{ij}|\mathcal E_{ij}|$ y $\operatorname{tr}W_c$
 (Secciones~\ref{sec:B4T01} y~\ref{sec:B4T04}) &
 Sección~\ref{sec:B5T03}, diferencia de Schur completa\\
Cambios de $\sum p_i^3$, $\sum p_i^2\kappa_i$ y $\sum_am_{2,a}^2$ (Sección~\ref{sec:B4T07}) &
 Sección~\ref{sec:B5T04}, prefactor y resolvente por entradas\\
Término cuadrático y resto cúbico del determinante (Sección~\ref{sec:B4T08}) &
 Sección~\ref{sec:B5T04}, cambio de soporte\\
\bottomrule
\end{tabular}
\end{center}

\par\emph{Cambio exacto de soporte y gauge.}
Sean $\Gamma=S_0^+$, $v=(\mathbf1_R,-\mathbf1_C,0_D)$,
$\Pi=vv^T/(2k)$ y $\widetilde S_0=S_0+\Pi$.
Sean $J$ las cuatro líneas retiradas e $I=J^c$. Ponemos
\[
 C_J=S_{0,JJ}+\Pi_{JJ},\quad V_J=S_{0,IJ}+\Pi_{IJ},
 \quad P_*=(\widetilde S_0)_{II}-V_JC_J^{-1}V_J^T,
\]
\[
 H_*=P_*^{-1}=(\Gamma+\Pi)_{II},\quad
 D_*=\Delta_0+(\Pi'-\Pi_{II})+V_JC_J^{-1}V_J^T.
\]
La identidad exacta de cofactores es
\[
 \frac{\det S'_{0,\rm gf}}{\det S_{0,\rm gf}}
 =\frac{k}{k-1}\frac{\det(I+H_*D_*)}{\det C_J}.
\]
La Sección~\ref{sec:B5T04} controla sus términos cuadrático y cúbico. No se usa ningún desarrollo
de Taylor hasta $p=0$ o $1$.

\par\emph{Conservación de la cancelación, con sus coeficientes.}
Sea $a_i$ el indicador de las cuatro líneas de $i$. Definimos
$f_i=\sum_{a<b\in i}\Gamma_{ab}$ y, con sólo los pares retenidos, $f_i^I$.
La respuesta cuadrática del hijo satisface
\[
 y_i/p_i=-a_i^T\Gamma a_c+O(k^{-2}),\qquad
 \mathcal Q(y)=a_c^T\Gamma a_c-4+12p_c+O(k^{-2}).
\]
La derivada exacta de Jacobi, y no la derivación de un resto asintótico, da
\[
 \delta_c=-6p_c-f_c+\sum_Fp_i(6p_i+f_i)
             -\sum_{E'}y_i(6p_i+f_i)+O(k^{-2}).
\]
La identidad de frontera de la acción, el cambio del prefactor y el cofactor dan
\[
 \begin{aligned}
 \epsilon_c:={}&L(B+c)-L(B)-\log p_c\\
 ={}&-6p_c-f_c+6\sum_Fp_i^2-6\sum_{E'}p_iy_i-\sum_{E'}y_if_i\\
 &+\sum_Fp_if_i^I-\sum_{j\in J}(\Gamma_{jj}-1)-1/k+O(k^{-2}).
 \end{aligned}
\]
La identidad final de la Sección~\ref{sec:B4T08}, antes de aproximar, implica
\[
 \sum_Fp_i(f_i-f_i^I)
 =-\sum_{j\in J}(\Gamma_{jj}-1)-1/k-p_cf_c.
\]
Como $p_cf_c=O(k^{-2})$, resulta $\epsilon_c-\delta_c=O(k^{-2})$.
Todas las sumas de errores usan $\sum_Ep=k$, $\sum_Fp=O(1)$ y la Sección~\ref{sec:B5T04}.
Este paso retiene y cancela los términos de orden $k^{-1}$.

La identidad de conservación de la fila da exactamente
\[
 \sum_{c\in r}p_c\delta_c=0,\quad
 w_c=p_c(1+\delta_c)>0,\quad\sum_{c\in r}w_c=1.
\]
La Sección~\ref{sec:B5T02} implica $\delta_c=O(k^{-1})$. Por tanto
\[
 r_c=\epsilon_c-\log(1+\delta_c)=O(k^{-2}),\quad
 \rho=\mathbb E_w e^{r_c},\quad|\log\rho|\le\sup_c|r_c|.
\]
Con esto queda probada la extensión determinista, bajo las hipótesis del teorema.
\end{proof}

\paragraph{Alcance.}
El caso de torres con máscara queda incluido mediante la representación con dos diagonales
de una sola casilla por casilla de la Sección~\ref{sec:B5T01}. La herencia es de un padre interior a una
clase
exterior con margen; no se afirma que una trayectoria permanezca en ella. No exhibimos aquí ningún
residual de reinas de esta clase con $k=o(n)$; el Corolario~\ref{cor:completo} sólo muestra que la
clase macroscópica contiene al tablero completo.
